% Лекция 17. Источник: Тервер(Лекция)_202611112317_42212.pdf, с. 2--4.
\clearpage
\begingroup
\everymath{\textstyle}
\subsubsection{Математическое ожидание произведения независимых величин}

\begin{theorem}
    Если $\xi$ и $\eta$ независимы и имеют конечные математические
    ожидания, то $\EE(\xi\eta)$ существует и
    \[
        \EE(\xi\eta)=\EE\xi\cdot\EE\eta.
    \]
\end{theorem}

\begin{Proof}[Дискретный случай]
    Пусть $\Pr\{\xi=x_i\}=p_i$ и $\Pr\{\eta=y_j\}=q_j$.
    По независимости
    \[
        \Pr\{\xi=x_i,\eta=y_j\}=p_iq_j.
    \]
    Абсолютная сходимость следует из равенства
    \[
        \sum_{i,j}|x_iy_j|p_iq_j
        =\left(\sum_i|x_i|p_i\right)
         \left(\sum_j|y_j|q_j\right)<\infty.
    \]
    Поэтому
    \[
        \EE\xi\cdot\EE\eta
        =\left(\sum_i x_ip_i\right)\left(\sum_j y_jq_j\right)
        =\sum_{i,j}x_iy_jp_iq_j=\EE(\xi\eta).
    \]
\end{Proof}

\subsubsection{Неравенство Йенсена}

\begin{theorem}[Неравенство Йенсена]
    Пусть $g\colon\R\to\R$ --- выпуклая вниз функция,
    а $\xi$ имеет конечное математическое ожидание. Тогда
    \[
        g(\EE\xi)\leq\EE g(\xi).
    \]
    Правая часть может быть равна $+\infty$.
\end{theorem}

\begin{Proof}
    Для любого $x_0\in\R$ существует число $\ell(x_0)$, такое что
    \[
        g(x)\geq g(x_0)+(x-x_0)\ell(x_0),\qquad x\in\R.
    \]
    Положим $x_0=\EE\xi$ и $x=\xi$:
    \[
        g(\xi)\geq g(\EE\xi)+(\xi-\EE\xi)\ell(\EE\xi).
    \]
    По монотонности и линейности математического ожидания
    \[
        \EE g(\xi)\geq g(\EE\xi)
        +\ell(\EE\xi)\underbrace{\EE(\xi-\EE\xi)}_{=\,0}
        =g(\EE\xi).
    \]
\end{Proof}

\newpage
\subsubsection{Неравенство Ляпунова}

\begin{theorem}[Неравенство Ляпунова]
    Если $0<s<t$ и $\EE|\xi|^t<\infty$, то
    \[
        \bigl(\EE|\xi|^s\bigr)^{1/s}
        \leq\bigl(\EE|\xi|^t\bigr)^{1/t}.
    \]
\end{theorem}

\begin{Proof}
    Из $|\xi|^s\leq1+|\xi|^t$ следует $\EE|\xi|^s<\infty$.
    Функция $g(x)=|x|^{t/s}$ выпукла, поскольку $t/s>1$.
    Применяя неравенство Йенсена к $\xi'=|\xi|^s$, получаем
    \[
        \bigl(\EE|\xi|^s\bigr)^{t/s}
        =g(\EE\xi')\leq\EE g(\xi')
        =\EE\bigl(|\xi|^s\bigr)^{t/s}=\EE|\xi|^t.
    \]
    Возводя обе части в степень $1/t$, получаем требуемое неравенство.
\end{Proof}

\subsection{Моменты случайной величины}

\begin{definition}
    \textbf{Моментом порядка $k$} случайной величины $\xi$, где
    $k\in\N$, называется число
    \[
        \EE\xi^k,
    \]
    если это математическое ожидание конечно, то есть
    $\EE|\xi|^k<\infty$.
\end{definition}

\begin{consequence}[Существование моментов]
    Из неравенства Ляпунова следуют утверждения:
    \begin{enumerate}
        \item Если существует конечный момент порядка $k$, то существуют
        все моменты меньших натуральных порядков $i<k$.
        \item Если конечного момента порядка $k$ нет, то нет и конечных
        моментов натуральных порядков $i>k$.
    \end{enumerate}
\end{consequence}

\begin{definition}[Характеристическая функция]
    \textbf{Характеристической функцией} случайной величины $\xi$ называется
    \[
        \varphi_\xi(t)=\EE e^{it\xi},\qquad t\in\R,\quad i^2=-1.
    \]
\end{definition}
\endgroup
